Evaluating creative work requires the kind of domain expertise that is scarce, expensive,
and difficult to apply consistently at scale. This paper reports an unexpected finding
about structured evaluation: a rubric applied without expert identity context systematically
\emph{inverts} the quality hierarchy, producing higher scores for accessible work than for
ambitious work. Across a fifteen-puzzle corpus spanning three quality tiers, rubric-only
evaluation yields a pass rate of 80\% for Standard Hunt puzzles but only 40\% for
MIT/Elite-tier puzzles---because elite puzzles deliberately sacrifice Clarity and
Solvability for Elegance and Reading Reward, and an equal-weight rubric treats that
trade-off as a defect. Expert identity, encoded in rich reviewer profiles, reintroduces
the calibration that converts raw dimension scores into verdicts appropriate to the paradigm
under review.

We ground this investigation in puzzle hunt design, a collaborative creative practice with
well-defined quality criteria, a recognized expert community, and the property that answer
correctness is binary while quality is contested. We assemble panels of AI reviewers
instantiated from persistent, domain-specific profiles---structured markdown documents
capturing each simulated reviewer's design philosophy, evaluative lens, and named
evaluative frameworks---and compare eight conditions in a controlled ablation across 15
real published puzzles from MIT Mystery Hunt 2023, Huntinality, and Teammate Hunt.

The eight-condition ablation reveals two additional findings. First, a \emph{feedback
quality gradient}: richer profiles surface six analytical frameworks (\emph{a-ha economy},
\emph{load-bearing test}, \emph{world-model consistency}, \emph{social fabric},
\emph{perceptual-shift}, \emph{visual grammar}) that generic profiles and bare prompting
do not introduce; these frameworks appear in C6 reviews but in no input text provided to
the reviewers. Second, adding the \emph{Riven Standard} as a scored seventh dimension
(``the puzzle IS what the field does, not an obstacle overlaid on it'') achieves a
24.5 percentage-point tier-separation gap between elite and magazine-tier puzzles,
compared to 10.9 points for the equal-weight baseline.

All results are characterized as valid within the authors' constructed evaluation
framework; a human calibration study remains future work. We release the full profile
library, scoring rubric, and replication corpus to support adaptation to other
creative evaluation domains.
